\cotefile{84}

\scanpage{35}

\uline{Dém} OPS $A=O[u]/(u^2+bu+c)$, et un \unsure{endom.|autom.} $\Theta$ de la forme envisagée doit transformer $u$ en \deleted{\unsure{et}}

\[
u'=\bar u+\alpha.1_A \qquad \alpha\in\Gamma\underline O
\]

(puisque $u'\equiv\bar u\equiv-u\bmod\underline O.1_A$), avec

\[
\left\{\begin{aligned}
\operatorname{Tr}(u')&=\operatorname{Tr}(u) (=-b)\\
N(u')&=N(u)(=c)
\end{aligned}\right.
\]

--- les $\Theta$ en question correspondent biunivoquement aux $u'$ (sans condition sur $A$)$^*$.

% slanted
\margin{Sauf la condition pour $\Theta$ en outre \unsure{pour|que} \unsure{l'être|n'être|être} l'identité sur un \deleted{\unsure{ouv.}} \unsure{un vois.} de $X$, qu'il faut regarder à part, en un point $s\in X_0\ldots$}

En termes de $\alpha$, les conditions s'écrivent respectivement

\[
\left\{\begin{aligned}
2\alpha&=0 &&(\text{car }\operatorname{Tr}\alpha.1_A=2\alpha)\\
\alpha(\alpha-b)&=0
\end{aligned}\right.
\]

En les pts de $X-X_0$, \deleted{\unsure{le}} \inserted{$2\alpha=0$} implique $\alpha=0$.

Reste à regarder en un point $s$ de $X_0$. \deleted{\unsure{Je} \unsure{dis}} \inserted{je dis que $\alpha_s\in M_s=\operatorname{rad}\underline O_{X,s}$ --- sinon,} \deleted{par le corps résiduel $k(s)$, on aura}

\[
\deleted{\unsure{\text{que}|\text{donc}}\qquad\alpha(s)(\alpha(s)-b(s))=0,}
\]

\deleted{mais par hypothèse \deleted{$b(s)$} $A(s)$ étale i.e. $b(s)\ne 0$, donc on doit avoir ou bien $\alpha(s)=0$, ou bien $\alpha(s)=b(s)$,} $\alpha_s$ serait inversible, et on aurait, à cause de \inserted{la relation} $\alpha(\alpha-b)=0$, $\alpha_s-b_s=0$ i.e. $\alpha_s=b_s$, \deleted{\unsure{Ma|Mais}} i.e. $\alpha=b$ au voisinage de $s$, \deleted{donc} et comme $2\alpha=0$, on aurait $2b=0$, donc comme $b_s$ est \uline{inversible} (car $A$ étale en $s$, et $s\in X_0$) $2.\deleted{\unsure{\text{sera}|\text{serait}}}1_{X}$ serait nul au vois. de $s$ i.e. $X$ de car. $2$.

% archive sheet number: 34
\scanpage{36}

\unsure{L'opération|L'application|Définition} $\sigma_e$ s'appelle \uline{symétrie} \uline{p.r.} à la \uline{droite} $D$, comme $\sigma$ est un autom. de $A$ pour la structure quad. définie par la norme, $\sigma_e$ est un autom. de la structure quadratique de $E_e$ pour $Q$, a fortiori un autom. de la structure de similitude.

\uline{Prop} Supposons qu'en tout point de $X_0=V(2.1_{\underline O_X})$, i.e. tout pt de car. rés. $2$, la structure de similitude donnée sur $E$ soit non dég. (i.e $A$ soit étale sur $\underline O$ en ce point). Alors pour tte droite \deleted{tot}alement non isotrope $D\subset E$, $\sigma_D$ est l'\uline{unique} autom. $\Theta$ de la structure de similitude $\Sigma$ (ou encore, de la structure quadratique, quand $\Sigma$ est donnée par une forme quadratique $Q$) \deleted{telle} qui induise $\mathrm{id}_D$ sur $D$, $-\mathrm{id}$ sur $E/D$,$^*$ et qui \inserted{\unsure{ne|en}} soit \inserted{l'identité au voisin.} \deleted{\unsure{en} \unsure{les} points $s$ $\unsure{\in}S$} \deleted{\unsure{d'étale} de chaque pt $s$ de $X$}\deleted{--- est \unsure{induit},} \inserted{\unsure{en les|qu'en} points $s$ de $X-X_0$)}

% slanted
\margin{Ceci équivaut à $\det\Theta=-1$ (compte tenu que $\Theta|D=\mathrm{id}_D$)}

\uline{NB} Cette dernière condition \unsure{bien|n'est} \unsure{sûre.|nécessaire} est \unsure{auto}

\uline{Cor.} Soit $A$ algèbre quadratique sur $X$, supposons qu'en tt $x\in X_0$, $A$ soit étale en $x$. Alors l'involution canonique est le seul automorphisme \inserted{d'Alg} de $A$ qui induise $-\mathrm{id}.$ sur $A/\underline O.1_A$.

% archive sheet number: 35
\scanpage{37}

Alors

\[
f: A\unsure{\longrightarrow|\mathbin{-}} E,\qquad u\longmapsto u.e
\]

est un iso respectant les structures quadratiques

\[
Q(f(u))=Q(u.e)=N(u)\underbrace{Q(e)}_{=1}=N(u)
\]

% slanted
\margin{--- Plus généralement, $Q(u.x)=N(u)Q(x)$ $(u\in\Gamma A,\ x\in\Gamma E)$}

par quoi la structure d'alg. de $A$ est transformée en l'unique str. d'Alg \inserted{quadratique} sur $\underline O$, admettant $e$ pour section unité, et $Q$ comme fonction quadratique Norme. La symétrie cherchée $\sigma_D$ peut être définie, par transport de structure, comme déduite de l'inv. canonique $\sigma$ de $A$

\[
\sigma:A\longrightarrow A\qquad \sigma(u)=\operatorname{Tr}.u-u.
\]

Donc on aura, par définition (posant $\sigma_e=\sigma_{\underline Oe}=\sigma_D$)

\[
\sigma_{\deleted{\unsure{D}}e}(ue)=\bar u.e
\]

On va donner une description équivalente, par

\[
\sigma_e(x)=\varphi_Q(e,x)e-x
\]

où $\varphi_Q$ est la forme bil. \inserted{sym} associée à $Q$. Cette formule ne fait que traduire celle pour \deleted{$A$}, $\sigma(u)=\operatorname{Tr}(u)-u$, compte tenu que \unsure{$\{$}

\[
\operatorname{Tr}(u)=\varphi_N(u,1).
\]

Avec la notation du ``produit scalaire'' (quand $2$ est inv sur $X$), on aura $\varphi_Q(e,x)=2e.x$, et on peut écrire

\[
\sigma_e(\deleted{\unsure{u}}x\deleted{\unsure{.e}})=2(e.x)e-x
\]

% archive sheet number: 36
\scanpage{38}

\deleted{Je dis \unsure{q}} Le torseur en question provient d'un torseur \inserted{$\Pi$} sous \deleted{\unsure{$\mathscr U_A$}} $U_A\subset A^*$, sous-schém. en groupes des sections de norme $1$, en posant

\[
\Pi=\left\{\Theta\in\Gamma\underbrace{SM^{*}}_{\substack{\text{sections}\\\text{de}\\\text{Norme }1\\(N(\Theta)=1)}}\ \middle|\ \begin{aligned}&\Theta\,\deleted{\unsure{u\mapsto u}}\ \text{normalise }A\\&\text{et y induit }u\longmapsto\bar u\end{aligned}\right\}
\]

4.2. Relation avec les symétries.

Soit $(E,\Delta',Q)$, $Q:E\longrightarrow\Delta'$, une structure de similitude str. fidèle. Soit $D$ un sous-fibré en droites de $E$, \deleted{\unsure{tot|tot.}} non isotrope i.e $Q|D$ est non nul sur chaque fibre. Je vais définir un automorphisme

\[
\sigma_D:E\xrightarrow{\sim}E
\]

de $E$ (avec sa structure de similitude), induisant $\mathrm{id}_D$ sur $D$, $-\mathrm{id}$ sur $E/D$. Quitte à tensoriser $E$ par $D^{-1}$, OPS $D$ muni d'une section base $e$, et profiter de celle-ci pour prendre une base $Q(e)$ dans $\Delta'$, donc OPS que $Q$ est \deleted{\unsure{définie par|défini par}} une fonction quadratique scalaire

\[
Q:E\mathbin{-}\underline O
\]

telle que

\[
Q(e)=1,
\]

% archive sheet number: 37
\scanpage{39}

au voisinage de $s$. Donc, en ce voisinage

\[
u'=\underset{\substack{\Vert\\-b-u}}{\bar u}+b=\deleted{(\operatorname{Tr}(u)1-u)}-u=u
\]

et donc on aurait $\Theta=\mathrm{id}$ au voisinage de $s$, contrairement à l'hypothèse.

Donc on conclut $\alpha_s\in M_s$, mais alors, comme $b_s\notin M_s$, \deleted{\unsure{$b_s\in O_s^*$}} on aurait $\alpha_s-b_s\notin M_s$ i.e. $\alpha_s-b_s\in\underline O_s^*$, d'où (comme $\alpha_s(\alpha_s-b_s)=0$) $\alpha_s=0$, donc $\alpha$ nul au voisinage de $s$, cqfd.

\uline{NB} En car. $2$, dans le cas non étale, l'exemple typique est $A=\underline O[u]/(u^2)$. Les autom. de $A$ sont donnés \inserted{(indép. de la car.)} par

\[
\Theta_{\lambda,\alpha}:u\longmapsto\lambda u+\alpha\qquad\lambda\in\Gamma O^*,\ \alpha\in\Gamma\underline O
\]

avec

\[
\left\{\begin{aligned}2\alpha&=0\\\alpha^2&=0\end{aligned}\right.
\]

ce qui, en car. $2$, se réduit à

\[
\alpha^2=0.
\]

Ceux qui induisent $-\mathrm{id}_{A/O.1}$ sont ceux de la forme

\[
u\longmapsto-u+\alpha\qquad(\text{avec }2\alpha=0,\ \text{\uuline{$\alpha^2=0$}}),
\]

donc en car. $2$ \unsure{ceux} \uline{de la forme} $u\longmapsto u+\alpha$ avec $\alpha^2=0$.$^*$

% slanted
\margin{Dire qu'en car $2$, \deleted{\unsure{les} \unsure{tot.}} ramifié ($\operatorname{Tr}=0$ \unsure{i.e.} $\sigma=\unsure{0|u}$) les symétries $\sigma_D$ sont l'\uline{identité}\,!}

% slanted
\margin{Il n'y a pas que l'identité i.e. pour que $\sigma_D$ puisse être l'identité sur $D$, $-\mathrm{id}$ sur $E/D$ (en car. $2$, les \unsure{car.} \unsure{réunies})}

Soit $(E,\Delta,Q)$ une str. de similitude str. fid, et $D\subset E$ droite tot\textsuperscript{t} non isotrope, d'où l'autom $\sigma_D\in\Gamma\underline{\operatorname{End}}(E)^*$. Le fait que $\sigma_D$ soit un autom de la structure de similitude signifie

% archive sheet number: 38
\scanpage{40}

simplement que

\[
\underbrace{\mathrm{int}(\sigma_D)}_{\substack{\text{autom de}\\\underline{\operatorname{End}}(E)\\\Theta\longmapsto\sigma_D\Theta\sigma_D^{-1}}}\ \text{transforme }\underset{\substack{\cap\\\underline{\operatorname{End}}(E)}}{A}\text{ en }\text{lui-}\mathord{=},
\]

et je dis qu'on a, plus précisément

\[
\boxed{\sigma_D u\sigma_D^{-1}=\bar u}\qquad\text{pour }u\in\Gamma A
\]

Pour le prouver, OPS $E=A$, \inserted{$\Delta=\underline O$, $Q=N_{A/\underline O}$,} $D=\underline O.1$, et on est ramené à prouver que

\[
\sigma\mu_u\sigma^{-1}=\mu_{\bar u}
\]

où $\mu_u\in\Gamma\underline{\operatorname{End}}(A)$ est défini par $\mu_u x=ux$.

La formule équiv. est alors :

\[
\overline{u\bar x}=\bar u x\qquad\text{OK.}
\]

\uline{Proposition} $(E,\Delta,Q)$, ou encore $(E,A\subset\operatorname{End}(E))$ structure de similitude str. fid. Soit

\[
\begin{aligned}
\Pi&\subset W(\underline{\operatorname{End}}(E))^*\\
\Vert&\\[-4pt]
\left\{\Theta\ \substack{\text{section loc}\\\text{de }\underline{\operatorname{End}}(E)\\\text{de }\mathrm{Sch}/X}\ \middle|\ \substack{\Theta\ \text{normalise }A,\ \text{induit }\sigma\text{ sur }A\\\text{et }\det\Theta=-1}\right\}&
\end{aligned}
\]

\deleted{Alors} \deleted{$\unsure{\Theta}$} et soit

\[
\begin{aligned}
M&\subset\unsure{W|\mathbb W}(\underline{\operatorname{End}}(E))\\
\Vert&\\[-4pt]
\{\Theta\mid\Theta u=\bar u\,\Theta\ \text{pour tte section }u\ \text{de }A\}&\\
&\quad\text{(ou pour \uline{une} section qui}\\[-2pt]
&\quad\text{engendre }A/\underline O.1_A\ldots)
\end{aligned}
\]

% archive sheet number: 39
\scanpage{41}

Alors

a) $M$ est un sous-fibré vectoriel de rang $2$ de $\underline{\operatorname{End}}(E)$, qui est un sous-$A$-Module à gauche et à droite. Les deux structures gauche et droite sont reliées par

\[
u\Theta=\Theta\bar u,\qquad\text{pour }\Theta\in\Gamma M,\ u\in\Gamma A
\]

(elles se déterminent \inserted{donc} mutuellement \ldots) Le $A$-Module $M$ à gauche \deleted{est loc libre} inv. (donc aussi celui à droite).

b) $\Pi$ est formé des sections $\Theta$ de $M$ telles que $N(\Theta)=-1$. Il forme un torseur$^*$ sous $SW(A)\subset W(A)^*$, le \deleted{groupe} schéma en groupes des $x\in\Gamma A$ tels que $N(x)=1$, par l'opération de multiplication à g. (\unsure{aussi bien|ou à droite} \inserted{par celle à} droite).

% slanted
\margin{\uline{NB} Ce torseur est en fait loc. trivial (Zar)}

\deleted{\uline{Dém.} OPS $E=A=\underline O[u]/(u^2+bu+c)$. Alors $W$ est défini comme le noyau de}

\[
\deleted{\unsure{\mathscr L_u}:\underline{\operatorname{End}}(\unsure{.}}
\]

c) Considérons

\[
\underset{\substack{\cup\\X}}{P}=P(E),\qquad P^*=P-\underset{\substack{\Vert\\\operatorname{Spec}(A)}}{X'}
\]

le fibré des droites tot. non isotropes de $A$.

$W(A)^*$ opère sur $P^*$ par son opération sur $E$, par $u\in\Gamma W(A)^*$ définit la transformation $D\longmapsto u.D$ (NB. $\unsure{G_m^{*}|G_m}=W(\underline O)^*\subset A^*$ opère trivialement donc c'est le quotient $W(A)^*/G_m$ qui opère \ldots).

On a un \deleted{\unsure{applicat|application}} morphisme canonique

\diagram{\begin{tikzcd}[ampersand replacement=\&,row sep=1pt]P^*\arrow[r,"f"]\&\Pi\\ D\arrow[r,maps to]\&\sigma_D\end{tikzcd}}

% archive sheet number: 40
\scanpage{42}

% vertical rule beside opening assertion
et on a

\[
f(u.D)=u\bar u^{-1}f(D)\ \inserted{\left(=f(D)\bar u^2\right)}
\qquad D\in\Gamma P^{\unsure{*|o}},\quad u\in\Gamma A^*
\]

i.e. le morphisme de schém. est compatible avec les opérations des groupes d'op. \inserted{$W(A)^*$ et $SW(A)^{\unsure{*}}$} et l'homomorphisme de groupes

\diagram{\begin{tikzcd}[ampersand replacement=\&,row sep=2pt]W(A)^*\arrow[r]\&SW(A)\\u\arrow[r,maps to,no head]\&u\bar u^{-1},\end{tikzcd}}

donc aussi (en restreignant le groupe d'opérateurs sur $P^*$ à $SW(A)$) avec

\[
u\longmapsto u^2\qquad SW(A)\longrightarrow SW(A)
\]

% archive sheet number: 41
\scanpage{43}

\mspage{(79)}

$\chi$-scindages, produits contractés d'extensions $\chi$-scindées, \deleted{produits} \unsure{construction|constructiō|contractés} \unsure{de revêtements quadratiques|du revêtement quadratique}, \unsure{invariants|invariant} d'Arf-Stiefel-Whitney généralisés, etc.

% stamp: MidiGraf Montpellier 2 1999 (vertical); geometric mark at lower right.

% archive sheet number: 42
\scanpage{44}

\uline{Revêtements quadratiques}

\diagram{\begin{tikzcd}[ampersand replacement=\&,column sep=small]0\arrow[r]\&\underline O\arrow[r]\& A\arrow[r,"\varphi"]\&\underset{\mathrm{inv}}{\check L}\arrow[r]\&0\end{tikzcd}\quad $E=\check A$\quad\begin{tikzcd}[ampersand replacement=\&,column sep=small]0\arrow[r,dash]\&L\arrow[r]\&E\arrow[r,"\psi"]\&\underline O\arrow[r]\&0\end{tikzcd}}

\[
\begin{aligned}
\underline P&=\psi^{-1}(1)\quad\text{tors. sous }L &\psi(x)&=\langle1,x\rangle\\
\underline P_\lambda&=\psi^{-1}(\lambda)\quad\text{id}&&=x(1)
\end{aligned}
\]

$\bullet\quad T\in\Gamma\underline P_2\subset E$\qquad fonction trace

$\bullet\quad N\in\Gamma\operatorname{Sym}^2(E)\subset\operatorname{Sym}^{\unsure{\bullet|*}}(E)$\qquad forme quadratique \uline{norme} (fonctions pol. sur $A$)

NB\qquad 1)

\diagram{\begin{tikzcd}[ampersand replacement=\&]0\arrow[r]\&E\otimes L\arrow[r]\&\operatorname{Sym}^2(E)\arrow[r,"\chi"]\&\underline O\arrow[r]\&0\end{tikzcd}}

\[
\begin{aligned}
&\text{où }\chi(Q)=Q(1)\\
&\text{on a }N\in\chi^{-1}(1)\quad(\text{tors. sous }E\otimes L)\\
&\text{i.e. }N(1)=1
\end{aligned}
\]

\deleted{\unsure{b)}} 2) \deleted{\unsure{On}} On a

\diagram{\begin{tikzcd}[ampersand replacement=\&]0\arrow[r,dash]\& L^{\otimes2}\arrow[r]\&E\otimes L\arrow[r]\&L\arrow[r]\&0\end{tikzcd}}

et $\operatorname{Sym}^2(E)$ a suite de comp. à quotients successifs $\underline O$, $L$, $L^{\otimes2}$. On identifie $L^{\otimes2}$ à ss-mod de $\operatorname{Sym}^2(E)$

\margin{attention au signe $\longrightarrow$ $\delta$ est le discr. de $N$ \uline{changé de signe}}

$\bullet\quad-\delta\in\Gamma(L^{\otimes2})$\qquad ``discriminant'' de $N$

\uline{NB}\qquad Forme bil. sym $\varphi_N$ associée à $N$

\[
\begin{aligned}
\varphi_N:\quad A\times A&\mathbin{-}\underline O\qquad\text{soit}\qquad A\longrightarrow\check A\overset{\mathrm{déf}}{\simeq}E\\
\text{d'où}\qquad \underset{\substack{\wr\\\check L}}{\bigwedge^2 A}&\longrightarrow\underset{\substack{\wr\\L}}{\bigwedge^2 E}\qquad\text{d'où }\delta\in\Gamma\underbrace{L^{\otimes2}}_{\substack{\wr\\\underline{\operatorname{Hom}}(\check L,L)}}
\end{aligned}
\]

% slanted
\margin{\uline{NB} On a $\varphi_N(x,y)=x\bar y+y\bar x$, et $\operatorname{Tr}(x)\operatorname{Tr}(y)=\varphi_N(x,y)+\operatorname{Tr}(xy)$ \unsure{ou i.e.|on trouve} $\varphi_N(x,y)=\operatorname{Tr}(x)\operatorname{Tr}(y)-\operatorname{Tr}(xy)$.}

\[
\boxed{T^2-4N=\inserted{+}\delta}\quad\text{i.e}\quad\boxed{4N=T^2\deleted{\unsure{+}}{\unsure{-}}\delta}.
\]

\margin{$\uparrow$ Attention au signe\,!}

\uline{NB} Soient $x',x''\in\Gamma E$ deux sections \uline{distinctes} en tt pt, et qui correspondent à des homs d'alg. $A\overset{x',x''}{\rightrightarrows}\underline O_X$ (donc $x',x''\in\Gamma\underline P_1$). On a donc $A\simeq\underline O_X\times\underline O_X$

\[
\begin{aligned}
\deleted{\unsure{x'(x+y)|x'(a,b)}}\qquad x'(a,b)&=a,\qquad x''(a,b)=b\\
\operatorname{Tr}(a,b)&=a+b,\qquad N(a,b)=ab
\end{aligned}
\]

d'où

\[
\boxed{\begin{aligned}T&=x'+x''\\N&=x'x''\\\delta&=(x'-x'')^2\end{aligned}}\qquad\deleted{\unsure{\delta=(x'-x'')^2}}
\]

% archive sheet number: 43
\scanpage{45}

% Blank verso. The faint reversed marks show through from the recto.
\scanpage{46}

Considérons la fonction polynomiale sur $\underline P_1$

\[
x\longmapsto (\underbrace{x-x'}_{\substack{\rotatebox{90}{$\scriptstyle\ni$}\\ L}})(\underbrace{x-x''}_{\substack{\rotatebox{90}{$\scriptstyle\ni$}\\ L}})\in L^{\otimes 2}
\]

\[
\text{i.e.}\qquad x\longmapsto x^2-Tx\deleted{\unsure{+N|\psi(x)}}+N
\]

restriction de la forme quadratique

\[
Q:\quad x\longmapsto x^2-\inserted{\psi(x)}Tx\deleted{\unsure{Q(x)|\psi(x)}}+\deleted{\unsure{N}}\psi(x)^2N,\qquad E\longrightarrow L^{\otimes 2}\subset\operatorname{Sym}^2(E)
\]

sur $E$. Cette forme est définie \inserted{par cette formule} en tous les cas, \unsure{supposée|supposons|sans supposer} existence de $x',x''$), elle prend ses valeurs dans $\operatorname{Sym}^2(E)$, mais \deleted{on} vérifie qu'elle est à valeurs dans $L^{\otimes 2}\subset\operatorname{Sym}^2$ (argument ``universel''\ldots) Sa restriction à $L$ est

\[
x\longmapsto Q(x)=x^2:\quad L\longrightarrow L^{\otimes 2}
\]

(partie ``hom.'' de la fonction polynomiale \inserted{de degré 2} envisagée sur $\underline P_1$).

\uline{NB} Equiv. de cat. fibrées

\[
\begin{aligned}
\substack{\text{rev. quadratiques}\\\text{de } S}
&\simeq\left\{\begin{aligned}
&\deleted{\text{Triples}}\ \text{Triples }(\deleted{\unsure{\underline P_1}}\underline L,\underline P_1,Q)\text{ sur }S\\
&\text{a) }\underline L\ \text{Module inv.}\\
&\text{b) }\underline P_1\quad \underline L\text{-torseur}\\
&\text{c) }Q:P_1\longrightarrow L^{\otimes 2}\quad\text{fonct. pol. de degré }2\\
&\text{à }\deleted{\text{partie}}\text{ homogène ``}x\longmapsto x^2\text{''}:\underline L\longrightarrow\underline L^{\otimes 2}
\end{aligned}\right.
\end{aligned}
\]

($Q$ est section d'un torseur sous $\deleted{A}\underline{\operatorname{Hom}}(\underset{\substack{\text{module envelop-}\\\text{pant de }P_1}}{E},\underline L^{\otimes 2})\simeq\mathscr A\otimes\underline L^{\otimes 2}$~; ext. de $\underline L$ par $\underline L^{\otimes 2}$~* \ldots).

% slanted
\margin{* Si $T$ $(=x'+x'')$ est donné d'avance (dans $\Gamma\underline P_2$), alors l'indétermination est \uline{seulement} du $N$, de façon que $x^2+Tx+N$ soit $\in\Gamma\underline L^{\otimes 2}$, \inserted{\unsure{donc indéterminé}} \unsure{donc} mod $\Gamma L^{\otimes 2}$. Si $2$ inv. sur $\mathscr O_S$, alors cette indétermination est levée par la \unsure{considération|donnée} de $\delta\in\Gamma\underline L^{\otimes 2}$, \unsure{vu la} \unsure{form.}~: $4N=T^2-\delta$ Mais il \deleted{\unsure{faut}} faut $T^2\equiv\delta\ (4)$}

Le rev\textsuperscript{t} quadr. associé à ces données peut être interprété comme le ``schéma des zéros de $Q$'', interprété lui-$=$ comme section du faisceau structural du fibré affine $P_1$ correspondant à $\underline P_1$.

% archive sheet number: 45
\scanpage{47}

Re NB \qquad Considérons

\[
\deleted{\unsure{f}}\varphi_N:\mathscr A\longrightarrow\mathscr A^\vee=E\qquad\text{via forme bil. ass. à forme qu. }N
\]

\[
\text{d'où}\qquad \mathscr A\xrightarrow{\deleted{\unsure{f}}\varphi_N}E\xrightarrow{Q}L^{\otimes 2}
\]

On trouve

\[
\boxed{Q(\varphi_N(x))=\deleted{\unsure{N(x)}}N(x)\delta}
\]

\margin{Vérification à partir du cas décomposé,}

\deleted{Si $\delta$ inv. (cas étale)}

Conditions équivalentes

a) $X'$ rev. \uline{étale} de $X$

b) $\varphi_N:\mathscr A\longrightarrow\mathscr A^\vee=E$ \uline{iso}

c) $\delta\in\Gamma L^{\otimes 2}$ section \uline{inv.} i.e. ne s'annule en aucun pt

Alors $\delta$ base de $L^{\otimes 2}$, $\varphi_N:\mathscr A\simeq E$, et en faisant les identifications, $Q$ correspond à $-N$

Si par contre $\delta=0$ (cas \deleted{\unsure{tott|tot\textsuperscript{t}}} ramifié) on trouve

\[
Q(\varphi_N(x))=0\ ..
\]

Considérons

\[
x\longmapsto\boxed{2x-T=z}\quad\inserted{(z\in\Gamma L\deleted{\unsure{^{\otimes 2}}})}\qquad \underline P_1\longrightarrow\underline L
\]

c'est un iso si 2 inv. sur $X$ donc iso sur $X_2=X-V(2)$. Si $x$ et $z$ se correspondent, on a

\[
2x=z+T\quad\Longrightarrow\quad Q(z+T)=4Q(x)
\]

et \deleted{le} \inserted{si $x$ section du} rev. quadratique $X'$ de $X$ i.e. \unsure{$\sqrt[2]{Q(x)=0}$|$Q(x)=0$}, on a $Q(T+z)=0$. Or on trouve \deleted{pour} \inserted{(comme $\psi(z+T)=2$)}

\[
\begin{aligned}
&\deleted{Q(T+z)}\\
&Q(z+T)=z^2\deleted{\unsure{+}} \unsure{-}\delta\\
&z^2+2Tz+\deleted{\unsure{T^2}}\unsure{-2}(Tz\unsure{-|+}\deleted{\unsure{T^2}})+4N=z^2\unsure{-}(T^2-4N)=z^2-\delta
\end{aligned}
\]

donc

\[
\left\{\begin{aligned}
&Q(x)=0\quad\text{i.e. }x\in\Gamma(X'/X)\quad\Longrightarrow\quad\boxed{z^2=\delta}\\
&(\text{et inv. si }2\text{ inv. sur }X.)
\end{aligned}\right.
\]

% archive sheet number: 46
\scanpage{48}

Le groupe $\mathbb Z/2\mathbb Z$ opère sur $X'/X$, via

\[
\left\{\begin{aligned}
&x\longmapsto\deleted{\unsure{T-x}}\bar x=T-x\quad\text{sur }\underline P_1\deleted{\unsure{^{*}}}\qquad(\text{auto})\\
&x+\bar x=T\qquad\text{si }x\in\Gamma\underline P_1
\end{aligned}\right.
\]

ou encore

\[
\left\{\begin{aligned}
&x\longmapsto\bar x=\psi(x)T-x\quad\text{sur }E\\
&\text{i.e.}\quad x+\bar x=\psi(x)T
\end{aligned}\right.
\]

On a

\[
\left\{\begin{aligned}
&\left.\begin{aligned}\psi(T)&=T\\\psi(N)&=N\end{aligned}\right\}\quad\nwarrow\\
&Q(\deleted{\unsure{\psi(x)}}\bar x)=Q(x)
\end{aligned}\right.
\]

Sur $\mathscr A$, l'involution donnée par

\[
x\longmapsto\bar x=\operatorname{Tr}(x)-x\qquad\text{i.e.}\qquad x+\bar x=\operatorname{Tr}(x).1
\]

Considérons $X_0=V(2)$, \deleted{\unsure{sur} \unsure{$T_0$} $\unsure{T\vert V(2)|N(x)}\in\Gamma\underline L$ et $T_0,N_0,$} \inserted{et $E_0,L_0,T_0,N_0,\unsure{Q_0}$ \inserted{$\delta_0$}} les restrictions à $X_0$. On a \deleted{$\unsure{T_0\in\Gamma|T_0\in E_0}$} alors

\[
T_0\in\Gamma\underline L_0\quad(\Gamma E_0)
\]

\[
\text{et}\qquad T_0^2=\delta_0\quad(=-\delta_0)\ \text{sur }X_0
\]

et $=$ mieux~: si $T_0$ est relevé loc\textsuperscript{t}. \inserted{en section $\tau_1$ de $E$} sur $X_1=V(4)$, \deleted{\unsure{on aura}} (p. ex. $\tau_1=T|X_1\overset{\mathrm{déf}}=T_1$) on aura

\[
\tau_1^2=\delta_1\quad(\overset{\mathrm{déf}}=\delta|X_{\unsure{4|1}})
\]

(NB cela ne dépend pas du choix de ce relèvement, car $(\tau_1+2\alpha)^2\equiv\tau_1^2\ (4)$) On l'exprimera par l'écriture

\[
\boxed{T_0^2\equiv\deleted{\unsure{-}}\delta\quad(4)}
\]

Ceci posé, si $x$ section de $X'/X$, $z$ la section corr. de $\underline L$, via $z=2x-T$, \deleted{\unsure{$x_0$|$z_0$}}, $z_0$ les restr. à $X_0$, on aura

\[
\left\{\begin{aligned}z_0&=T_0\\z^2&=\delta\end{aligned}\right.
\]

i.e $z$ est une solution de $z^2=\delta$ \inserted{sur $X$} qui prolonge la solution déjà donnée \unsure{$z_0$|$T_0$} sur $X_0$.

% archive sheet number: 47
\scanpage{49}

Soit

\[
\unsure{\mathcal X'|\mathfrak X'}:\mathrm{Sch}_{/X}\longrightarrow\mathrm{Ens}\qquad\text{dont la val. sur }X\text{ soit}
\]

\[
\mathcal X'(X)=\left\{x\in\Gamma\underline L\ \middle|\ \begin{aligned}x^2&=\delta\\x|X_0&=T_0\end{aligned}\right\}\qquad\begin{gathered}\text{et idem via}\\Y\longrightarrow X.\end{gathered}
\]

\[
\deleted{\unsure{\text{Soit}}}\quad\mu'_{2\unsure{X|x}}\subset\mu_{2X}:\mathrm{Sch}_{/X}\longrightarrow\text{Groupes comm}
\]

\[
\mu'_2\unsure{=}(X)=\{\lambda\in\Gamma\underline{\mathscr O}_X^*\mid\lambda^2=1,\ \lambda_0=1\}
\]

% slanted
\margin{\uline{NB} $\mathcal X'$ et $\mu'_2$ sont des faisceaux fpqc, \unsure{$\varnothing$|$\mathscr O$} \uline{pas} repr. \deleted{\unsure{si}} $\exists\,x\in X$ tel que, \deleted{\unsure{soit}} \deleted{\unsure{nul soit}} $2$ soit ni inv. ni nul ds $\mathscr O_{X,x}$, i.e. $2\in\mathfrak m_{X,x}$, $2\ne0$}

Alors $\mu'_2$ opère sur $\mathcal X'$, et si $\delta$ est inv. i.e. si $X'/X$ étale, alors $\mathcal X'$ est un torseur sous $\mu'_2$ (alors que $X'/X$ est un torseur sous $(\mathbb Z/2\mathbb Z)_X$).

On a compatibilité

\diagram{\begin{tikzcd}[ampersand replacement=\&]
X' \arrow[r,"f"] \& \mathcal X' \\
(\mathbb Z/2\mathbb Z)_X \arrow[r] \& \mu'_{2X}
\end{tikzcd}}

i.e.

\[
f(\bar x)=-\deleted{\unsure{\bar x}}f(x)
\]

où

\[
z\longmapsto-z:\quad\mathcal X'\longrightarrow\mathcal X'
\]

est l'\deleted{application}\unsure{automorphisme|opération} induit par la section $-1$ de $\mu'_2$ sur $X$ \deleted{\illegible{2}}.

\smallskip\centerline{\rule{1em}{0.3pt}}\smallskip

Supposons $2$ $\underline{\mathscr O}_X$-régulier. Alors la donnée des rev. quadratiques $X'/X$ \uline{équivaut} à celle des

\[
\left\{\begin{aligned}
&\underset{\text{inv.}}{L},\quad\delta\in\Gamma L^{\otimes 2},\quad T_0\in\Gamma(L_0)\qquad(\text{où }L_0=L|X_0)\\
&\text{satisfaisant}\qquad T_0^2\equiv\delta\quad(4)
\end{aligned}\right.
\]

On veut décrire les sections $x$ de $X'/X$, en termes de la section $z$ de $L$, dans le cas où $X'\longrightarrow\mathcal X'$ pas iso i.e. où $X_0\ne\varnothing$. \deleted{On aura} \inserted{Bien sûr,} on aura

% archive sheet number: 48
\scanpage{50}

\[
\Gamma(X'/X)\lhook\joinrel\longrightarrow\Gamma(\unsure{\mathcal X'|\mathfrak X'}/X)\qquad(\text{injectif})
\]

car $X-X_0$ schém. dense dans $X$. Donc $x$ connue quand on connait $z$. Quelle condition à mettre sur $z$~? \deleted{Il faut que} Notons qu'on a

\[
\deleted{\unsure{T_0+z_0}}\qquad T+z\equiv0\quad(2)
\]

car cela signifie $T_0+z_0=0$ i.e. $z_0=T_0$, \inserted{OK}, donc on pourra prendre

\[
x=\frac{T+z}{2},
\]

je dis que a) c'est une section de $P_1$ \unsure{à} i.e. $\psi(x)=1$ (or cela équivaut à $2\psi(x)=2$ i.e. $\underset{\substack{''\\\psi(T+z)}}{\psi(2x)}=2$, Ok),

et b) on a bien $Q(x)=0$ (car ceci équiv. à $4Q(x)=0$ \unsure{i.e.|i.,} $Q(2x)=0$ i.e. $Q(T+z)=0$, donc à $z^2=\delta$ \ldots).

Description de $E$ en termes de $L,\delta,T_0$~: soit \inserted{(loc\textsuperscript{t})} $\tau$ une section de $L$ \inserted{sur $X$} qui remonte $T_0$

\[
\Gamma E=\left\{\underbrace{(\lambda,\unsure{\eta|y|\xi})}_{x}\ \middle|\ \lambda\in\Gamma\underline{\mathscr O},\ x\in\Gamma\left(\underbrace{\frac{\lambda\tau}{2}+L}_{\substack{\text{ne dépend}\\ \text{pas du choix}\\ \text{de }L}}\right)\right\}
\]

\[
\{\lambda=\psi(x),\ \eta=x-\psi(x)\tfrac{T}{2}
\]

\margin{\uline{NB} $\dfrac{\lambda\tau}{2}+L\subset\dfrac12\underline L$}

ou encore

\[
\boxed{\Gamma E\simeq\bigl\{(\lambda,\zeta)\mid\lambda\in\Gamma\underline{\mathscr O},\quad\zeta\in\Gamma(\underbrace{\lambda\tau+2L}_{\overset{\mathrm{déf}}=\ L_\lambda\text{ sous-}\deleted{\unsure{\text{Module}}}\text{torseur de }\underline L})\subset\Gamma(L)}
\]

* i.e. section de $L$ qui remonte la section $\lambda T_0$ de $L_0$ i.e. $\boxed{\zeta_0=\lambda_0T_0}$

\[
\boxed{\left\{\begin{aligned}\lambda&=\psi(x)\\\zeta&=2x-\psi(x)T\end{aligned}\right.}
\]

% slanted
\margin{* via l'op. de $\underline L$ dessus~: $U_\zeta x=x+\underline{2}\zeta$}

Relation entre $(\lambda,\eta)$ et $(\lambda\unsure{,}\zeta)$~: $\zeta=2\eta$

\margin{\deleted{\uline{NB} \unsure{On a} $\boxed{\left\{\begin{aligned}\lambda&=\psi(x)\\\zeta&=2x-\psi(x)T\end{aligned}\right.}$ pour $x=(\lambda,\zeta)$}}

\[
0\longrightarrow\underline L\longrightarrow E\xrightarrow{\psi}\underline{\mathscr O}\longrightarrow0\qquad\left\{\begin{aligned}\psi(\lambda,\zeta)&=\lambda\\i(\unsure{\eta|y|\xi})&=(0,2\eta)\end{aligned}\right.
\]

$T$ correspond à la section $(2,\deleted{\unsure{\zeta}}0)$ de $\Gamma P_2\simeq\Gamma(\underbrace{2\tau+2L}_{2L})$

% archive sheet number: 49
